% Options for packages loaded elsewhere
\PassOptionsToPackage{unicode}{hyperref}
\PassOptionsToPackage{hyphens}{url}
\documentclass[
  ignorenonframetext,
]{beamer}
\newif\ifbibliography
\usepackage{pgfpages}
\setbeamertemplate{caption}[numbered]
\setbeamertemplate{caption label separator}{: }
\setbeamercolor{caption name}{fg=normal text.fg}
\beamertemplatenavigationsymbolsempty
% remove section numbering
\setbeamertemplate{part page}{
  \centering
  \begin{beamercolorbox}[sep=16pt,center]{part title}
    \usebeamerfont{part title}\insertpart\par
  \end{beamercolorbox}
}
\setbeamertemplate{section page}{
  \centering
  \begin{beamercolorbox}[sep=12pt,center]{section title}
    \usebeamerfont{section title}\insertsection\par
  \end{beamercolorbox}
}
\setbeamertemplate{subsection page}{
  \centering
  \begin{beamercolorbox}[sep=8pt,center]{subsection title}
    \usebeamerfont{subsection title}\insertsubsection\par
  \end{beamercolorbox}
}
% Prevent slide breaks in the middle of a paragraph
\widowpenalties 1 10000
\raggedbottom
\AtBeginPart{
  \frame{\partpage}
}
\AtBeginSection{
  \ifbibliography
  \else
    \frame{\sectionpage}
  \fi
}
\AtBeginSubsection{
  \frame{\subsectionpage}
}
\usepackage{iftex}
\ifPDFTeX
  \usepackage[T1]{fontenc}
  \usepackage[utf8]{inputenc}
  \usepackage{textcomp} % provide euro and other symbols
\else % if luatex or xetex
  \usepackage{unicode-math} % this also loads fontspec
  \defaultfontfeatures{Scale=MatchLowercase}
  \defaultfontfeatures[\rmfamily]{Ligatures=TeX,Scale=1}
\fi
\usepackage{lmodern}
\ifPDFTeX\else
  % xetex/luatex font selection
\fi
% Use upquote if available, for straight quotes in verbatim environments
\IfFileExists{upquote.sty}{\usepackage{upquote}}{}
\IfFileExists{microtype.sty}{% use microtype if available
  \usepackage[]{microtype}
  \UseMicrotypeSet[protrusion]{basicmath} % disable protrusion for tt fonts
}{}
\makeatletter
\@ifundefined{KOMAClassName}{% if non-KOMA class
  \IfFileExists{parskip.sty}{%
    \usepackage{parskip}
  }{% else
    \setlength{\parindent}{0pt}
    \setlength{\parskip}{6pt plus 2pt minus 1pt}}
}{% if KOMA class
  \KOMAoptions{parskip=half}}
\makeatother
\setlength{\emergencystretch}{3em} % prevent overfull lines
\providecommand{\tightlist}{%
  \setlength{\itemsep}{0pt}\setlength{\parskip}{0pt}}
\usepackage{bookmark}
\IfFileExists{xurl.sty}{\usepackage{xurl}}{} % add URL line breaks if available
\urlstyle{same}
\hypersetup{
  hidelinks,
  pdfcreator={LaTeX via pandoc}}

\author{\texorpdfstring{}{}}
\date{}

\begin{document}

\section{\texorpdfstring{Discussion: Implications and Community
Recommendations
\label{sec:discussion}}{Discussion: Implications and Community Recommendations }}\label{discussion-implications-and-community-recommendations}

\begin{frame}{Relationship to Prior Development Directions}
\protect\phantomsection\label{relationship-to-prior-development-directions}
Knight, Cordes, and Friedman \citep{knight2022fep} identified six
development directions for systematic Active Inference literature
analysis: (1) increased scope of relevant works, (2) richer annotation
schemes, (3) integration of manual and artificial contributions, (4)
transferable approaches across fields, (5) participation by diverse
contributors, and (6) updated analyses tracking the field's evolution.
This pipeline directly addresses directions 1, 2, 3, and 6: it scales
retrieval to three databases, replaces manual annotation with LLM-driven
extraction while preserving human review pathways, and produces a
pipeline designed for incremental re-execution as new literature
appears. Directions 4 and 5---cross-field transferability and community
participation---remain open and are addressed below.
\end{frame}

\begin{frame}{Tactical and Strategic Priorities}
\protect\phantomsection\label{tactical-and-strategic-priorities}
\begin{block}{Adopt Rigorous Reporting Metadata}
\protect\phantomsection\label{adopt-rigorous-reporting-metadata}
Papers should systematically report DOIs, ORCIDs, and explicit
hypothesis commitments. Submitted preprints should forward-link to their
published versions to prevent fragmented citation subgraphs. Our
extraction pipeline prioritizes the DOI as the canonical identifier;
failing that, deduplication cascades to arXiv IDs, Semantic Scholar IDs,
and OpenAlex IDs. Broad DOI adoption would resolve the cross-source
mismatch problem, enabling higher-resolution evidence mapping.
\end{block}

\begin{block}{Explore Open Knowledge Graph Infrastructure}
\protect\phantomsection\label{explore-open-knowledge-graph-infrastructure}
We encourage the exploration of federated nanopublication server
architectures to house community-contributed assertions. This would
enable a continuously updated living literature review that incorporates
new findings as they are published. The release of nanopub-js v0.1.0
\citep{kuhn2026nanopubjs} makes browser-based creation and querying of
nanopublications practical, enabling researchers to contribute
assertions directly from web interfaces. Integrating this approach with
the Active Inference Institute's Knowledge-Engineering infrastructure
\citep{knight2022fep} could provide the standardized semantic vocabulary
necessary for rigorous cross-study comparison.
\end{block}

\begin{block}{Standardize the Ontological Lexicon}
\protect\phantomsection\label{standardize-the-ontological-lexicon}
Immediate future extraction cycles should align assertion predicates
with the formally curated Active Inference Ontology. Enforcing shared
ontological primitives across studies will accelerate the aggregation of
evidence from otherwise siloed research communities, advancing the
interoperability goal outlined by Knight et al.~\citep{knight2022fep}.
\end{block}
\end{frame}

\begin{frame}{Empirical and Theoretical Imperatives}
\protect\phantomsection\label{empirical-and-theoretical-imperatives}
\begin{block}{Architect Unified Performance Benchmarks}
\protect\phantomsection\label{architect-unified-performance-benchmarks}
The computational tools domain (B) lacks standardized performance
benchmarks for direct comparison against deep reinforcement learning
architectures. Establishing baseline metrics analogous to standard RL
environments (e.g., OpenAI Gym) is a prerequisite for transitioning
theoretical proposals into applied systems.
\end{block}

\begin{block}{Prioritize Empirical Validation}
\protect\phantomsection\label{prioritize-empirical-validation}
Biology (C5) and Language (C3) have established theoretical frameworks
but limited empirical validation. Targeted experiments designed to test
specific FEP-derived predictions---such as demonstrating morphogenesis
as Bayesian inference or measuring active inference advantages in
language tasks---would strengthen the evidence base beyond what further
theoretical work alone can achieve.
\end{block}
\end{frame}

\begin{frame}[fragile]{Living Review Maintenance}
\protect\phantomsection\label{living-review-maintenance}
The pipeline is designed for continuous operation rather than one-time
analysis. Incremental resume capabilities (checkpoint-based assertion
extraction, merge-on-add corpus deduplication) enable periodic
re-execution as new papers are indexed. We envision a maintenance cycle
in which the pipeline is re-run quarterly, with updated hypothesis
scores and field statistics published alongside the pipeline release.
Community contributors can extend the framework by adding custom
hypothesis definitions, alternative keyword taxonomies, or
domain-specific extraction prompts---all configurable via the YAML
configuration file without modifying source code. A complementary
long-term trajectory is toward RAG-enabled access: integrating the
nanopublication knowledge graph into a Retrieval-Augmented Generation
architecture \citep{fan2024survey} would enable natural-language
querying of the evidence base, making quantitative literature synthesis
accessible to researchers without programming expertise.

\begin{block}{Agentic Workspaces and MCP Integration}
\protect\phantomsection\label{agentic-workspaces-and-mcp-integration}
Beyond traditional open-source maintenance, the repository is
architected as an intrinsically agentic workspace. Every underlying
source module (e.g., \texttt{src/knowledge\_graph/},
\texttt{src/visualization/}) is governed by dedicated \texttt{SKILL.md}
files serving as Model Context Protocol (MCP) prompt-boundaries. These
explicitly define the ``rules of engagement'' for autonomous AI
inference agents---such as enforcing the zero-mock testing philosophy
via local HTTP proxies, handling specific LLM fallback parsing logic,
and respecting headless rendering constraints. This design ensures that
future AI orchestrators can natively interface with, scale, and refine
the computational meta-analysis pipeline safely and deterministically
without structural micromanagement.
\end{block}
\end{frame}

\begin{frame}{Open Questions}
\protect\phantomsection\label{open-questions}
This meta-analysis surfaces four critical, empirically testable
questions warranting dedicated investigation:

\begin{itemize}
\tightlist
\item
  \textbf{Classifier calibration:} What proportion of A1 (Formal Theory)
  papers would be reclassified under an embedding-based or
  expert-annotated scheme, and how does this affect our understanding of
  the field's theoretical core?
\item
  \textbf{Falsifiability and Explicit Testing:} H1 (FEP Universality)
  produces a predominantly neutral evidence profile, consistent with the
  critique that FEP accommodates any behavior without generating
  distinctive predictions \citep{colombo2021free}. Can hypothesis
  definitions (and author reporting standards) be reformulated to
  require formal demonstration of a specific, refutable empirical
  prediction before contributing a supporting assertion?
\item
  \textbf{The Scalability Gap:} H5 (AIF Scalability) shows a strong
  positive trend, yet head-to-head comparisons with deep RL remain
  concentrated on a specific subset of benchmarks. Beyond what
  state-space dimensionality and reward density does the performance
  advantage of model-based AIF (via expected free energy exploration)
  degrade relative to model-free architectures?
\item
  \textbf{Evidence Cross-Pollination:} To what extent do mathematical
  structures underlying variational free energy minimization and energy
  function optimization in Energy-Based Models (e.g., VAEs, contrastive
  divergence) converge? Identifying shared architectural insights at
  this intersection could accelerate both Active Inference tools and
  mainstream machine learning.
\end{itemize}
\end{frame}

\end{document}
